% !PRINTED 625
\weberpage{625}
%<<CONT
congruent to a power of $p$ when $\lambda$ is divisible by $e$. For only then can the congruence
\[ \lambda \equiv \gamma x \, (\mathrm{mod}\ \mu) \]
be satisfied. If we set
\[ \gamma = e \nu, \]
then $\nu$ is relatively prime to $f$. From this it follows that the sequence of numbers
\weq{9}{1, g, g^2, \ldots, g^{e-1}}
can be taken for the $\xi$. For if $h, h'$ are two different numbers of the sequence $0, 1, 2, \ldots, e-1$, then $h'-h$ is never divisible by $e$, and consequently
\[ \xi' \xi^{-1} = g^{h'-h} \]
is never congruent to a power of $p$.

We shall further distinguish the two cases in which $p-1$ is or is not divisible by $q$.

a) If $p \equiv 1 \, (\mathrm{mod}\ q)$, then $\gamma \equiv 0 \, (\mathrm{mod}\ q-1)$, and consequently $e$ is divisible by $q-1$. We can therefore set
\[ e = \varphi(q^{\varkappa_1}) = q^{\varkappa_1 - 1}(q-1) \]
where $0 < \varkappa_1 \leqq \varkappa$.

When $\varkappa_1 = \varkappa$, and hence $e = \varphi(m)$, $f=1$, and therefore $p-1$ is divisible by $m$, then $\xi$ runs through the whole group $\mathfrak{N}$, and all the prime functionals $\mathfrak{p}_n$ are different from one another.

If $\varkappa_1 < \varkappa$, then $\nu$ is not divisible by $q$, since otherwise $e$ would not be the greatest common divisor of $\mu$ and $\gamma$, and
\[ p - 1 \equiv g^{\nu\varphi(q^{\varkappa_1})} - 1 \, (\mathrm{mod}\ m) \]
is divisible by $q^{\varkappa_1}$, but by no higher power of $q$. If we therefore set $q^{\varkappa_1} = m_1$ and
\weq{10}{m = m_1 m_2,}
then $m_1$ is the \gesperrt{greatest common divisor of $p-1$ and $m$}. The prime $p$ splits in the field $\Omega_m$ into $\varphi(m_1)$ different prime factors. But $p$ splits into just as many prime factors in the field $\Omega_{m_1}$, which is a divisor of the field $\Omega_m$ and consists of all rational functions of the $m_1$th root of unity $r_1 = r^{m_2}$. From this it follows that the prime factors of $p$ in $\Omega_{m_1}$ are not further decomposable in the field $\Omega_m$ either, and that we can therefore represent the prime factors $p$ in $\Omega_m$ as functionals of the field $\Omega_{m_1}$.
% [sic] Source prints 'Primfactoren p' without 'von'; retain the printed wording.

Here $\xi$ runs through a complete system of residues not divisible by $q$ with respect to the modulus $m_1$, and the $\mathfrak{p}_\xi$ are ideals in $\Omega_{m_1}$.
